\chapterhead{56}{ORR-17}{Stress Testing Banks}{4}{0}

\lohead{56}{a}{Describe the evolution of the stress testing process and compare
the methodologies of historical European Banking Association (EBA), Comprehensive
Capital Analysis and Review (CCAR), and Supervisory Capital Assessment Program
(SCAP) stress tests.}

\orsub{Supervisory Capital Assessment Program (SCAP)}

\noindent\begin{minipage}{\textwidth}
\renewcommand{\arraystretch}{1.32}
\begin{tabularx}{\textwidth}{@{}YY@{}}
\rowcolor{navy}\multicolumn{2}{@{}c@{}}{\thd{Comparison of stress testing pre-SCAP and post-SCAP}}\\
\rowcolor{palenavy} \tsub{Pre-SCAP} & \tsub{Post-SCAP}\\
{\tabfont Primarily assessed exposure to single-shocks (e.g., volatility
increases OR interest rate increases OR increasing unemployment).} &
{\tabfont Considers broad macro-scenarios and market-wide stresses with multiple
factors occurring or changing at once, as evidenced in the 2007--2009 financial
crisis.}\\
\rowcolor{paleblue}
{\tabfont Focused on specific bank products or business units (e.g., lending or
trust).} &
{\tabfont Focuses on the whole firm, a more comprehensive look at the effect of
the stress scenarios on the institution.}\\
{\tabfont Typically focused on earnings shocks (i.e., losses) but not on capital
adequacy.} &
{\tabfont Explicitly focuses on capital adequacy. Considers the post-stress
common equity threshold to ensure that a bank remains viable.}\\
\rowcolor{paleblue}
{\tabfont Focused exclusively on losses.} &
{\tabfont Focuses on revenues, costs, and projected losses.}\\
{\tabfont Stress testing was static in nature.} &
{\tabfont Stress testing is now dynamic and path dependent.}\\
\end{tabularx}
\end{minipage}
\orfigcap{What SCAP changed about stress testing.}

% ---------------------------------------------------------------------
\lohead{56}{b}{Explain challenges in designing stress test scenarios, including
the problem of coherence in modeling risk factors.}

\orsub{Challenges in designing stress tests}

Designing stress tests is difficult because they need to be both extreme and
realistic (coherent). This is because of:

\begin{lettered}
\item \textbf{Multi-factored problems.} Real-world problems involve many
interconnected factors. Changing one factor, like increased volatility, can
trigger changes in others, like credit markets freezing.
\item \textbf{Specifying joint outcomes.} A stress test needs to consider how
various risk factors will behave together. For example, high unemployment likely
coincides with falling equity prices.
\item \textbf{Not everything goes bad simultaneously.} Stress tests should account
for offsetting effects. If some currencies depreciate, others should appreciate.
During a ``flight to quality'', there must be safe-haven assets identified.
\item \textbf{Complexity in marked-to-market portfolios.} These portfolios involve
numerous positions and risk factors, making it difficult to model coherent
outcomes due to the sheer number of variables.
\end{lettered}

\orsubb{Addressing the challenges}

\begin{enumerate}
\item[\textcolor{rust}{\sffamily\bfseries 1.}] \textbf{Simple versus complex
scenarios.} The 2009 SCAP used simpler scenarios with a few key variables based on
historical data, achieving some level of coherence.
\item[\textcolor{rust}{\sffamily\bfseries 2.}] \textbf{Shift from
one-size-fits-all.} Previously, all banks faced identical stress tests. Now, banks
can submit their own scenarios alongside the supervisory ones to capture
bank-specific vulnerabilities. This allows regulators to gain valuable insights
into potential risks.
\end{enumerate}

% ---------------------------------------------------------------------
\lohead{56}{c}{Explain challenges in modeling a bank's revenues, losses, and its
balance sheet over a stress test horizon period.}

\orsub{Challenges in modeling losses and revenues}

Current stress tests rely on macro-factors (unemployment, GDP growth) but struggle
to translate these into micro, bank-specific outcomes like revenue and losses.
They struggle because of:

\begin{enumerate}
\item[\textcolor{rust}{\sffamily\bfseries 1.}] \textbf{Mapping macro to micro.}
It is difficult to map broad macro-economic factors to specific bank losses
because they affect different products and regions differently.\par
{\fontsize{9}{11}\selectfont\color{rust}\itshape E.g.\ credit card losses are
higher in areas with high unemployment, like Nevada in 2011.}
\item[\textcolor{rust}{\sffamily\bfseries 2.}] \textbf{Asset class
sensitivities.} Different asset classes react differently to economic
downturns.\par
{\fontsize{9}{11}\selectfont\color{rust}\itshape E.g.\ during a recession, fewer
cars are sold overall, but used car sales might increase as people seek
affordability.}
\end{enumerate}

\begin{ornotebox}{Modeling revenues versus losses}
Modelling revenue under stress is less developed than modelling losses. The 2009
SCAP offered little guidance on calculating stressed revenues.
\end{ornotebox}

\orsub{Challenges in modeling the balance sheet}

Bank stress testing involves a two-year projection of a bank's balance sheet and
income statement to assess whether it has enough capital to withstand economic
stress. This process entails:

\begin{checks}
\item Evaluating capital ratios using common equity as a base measure.
\item Calculating risk-weighted assets (RWAs) based on asset type and risk.
\item Making strategic decisions about asset sales, dividends, and capital
actions, which affect the balance sheet.
\item Maintaining appropriate capital and liquidity levels throughout the testing
period.
\item Projecting additional reserves for a year beyond the stress test to cover
future loan and lease losses.
\end{checks}

The challenge lies in the dynamic nature of the balance sheet over this period,
given the various strategic decisions banks may need to make in response to the
stress scenario.

\orsub{Stress test comparisons}

\noindent\begin{minipage}{\textwidth}
\renewcommand{\arraystretch}{1.32}
\begin{tabularx}{\textwidth}{@{}L{0.26\textwidth}Y@{}}
\rowcolor{navy} \thd{Test} & \thd{Disclosure and approach}\\
{\tabfont \textbf{1) 2009 SCAP} (U.S.)} &
{\tabfont Offered detailed disclosures including projected losses for each bank
across eight asset classes, and other resources to absorb losses. This high
transparency allowed market participants to evaluate stress test severity and
outcomes at the individual bank level.}\\
\rowcolor{paleblue}
{\tabfont \textbf{2) 2011 CCAR} (U.S.)} &
{\tabfont Required publication of only macro-scenario results, not specific
bank-level results, reducing transparency compared to the 2009 SCAP.}\\
{\tabfont \textbf{3) 2012 CCAR} (U.S.)} &
{\tabfont Returned to the disclosure level similar to the 2009 SCAP, with detailed
bank-level stress data including loss rates and losses by major asset classes.}\\
\rowcolor{paleblue}
{\tabfont \textbf{4) 2011 European stress tests}} &
{\tabfont Provided extensive detail post-CCAR, especially the Irish stress test,
which included a comparison of bank and third-party loss estimates. The
availability of data in a downloadable format aimed to restore credibility and
transparency after previous less effective tests.}\\
\end{tabularx}
\end{minipage}
\orfigcap{Four stress tests, compared on transparency.}
